\documentclass[14pt,aspectratio=169]{beamer}
\usepackage{PTSans}
\usepackage[T2A]{fontenc}
\usepackage[utf8]{inputenc}
\usepackage[russian]{babel}
\setbeamertemplate{navigation symbols}{}
\setbeamertemplate{footline}[frame number]
\setbeamerfont{frametitle}{size=\large}
\setbeamersize{text margin left=8mm,text margin right=8mm}
\begin{document}
\begin{frame}{Цель и условия}
    \textbf{Определение сопротивления резистора}
    \begin{itemize}
        \item Цель: оценить сопротивление по напряжению и току.
        \item Модель: $R=U/I$ для выбранного режима измерения.
        \item Условия: установившийся ток, температура считается постоянной.
    \end{itemize}
    \alert{Учебный пример: все данные вымышлены.}
\end{frame}

\begin{frame}{Установка и метод}
    \begin{itemize}
        \item Амперметр включен последовательно с резистором.
        \item Вольтметр подключен параллельно резистору.
        \item Снимаются пары $U$ и $I$ при разных напряжениях.
    \end{itemize}
    В этой упрощенной модели влияние приборов на цепь считается
    пренебрежимо малым. В реальном опыте это допущение нужно обосновать.
\end{frame}

\begin{frame}{Учебные данные}
    \begin{center}
        \begin{tabular}{c|rrrr}
            $U$, В & 1 & 2 & 3 & 5 \\
            \hline
            $I$, А & 0,020 & 0,041 & 0,059 & 0,100
        \end{tabular}
    \end{center}
    Для подробного расчета выбрана последняя точка.

    Таблица показывает исходные пары. Проверка линейности
    требует отдельного анализа всей серии.
\end{frame}

\begin{frame}{Расчет в выбранной точке}
    \[
        U=5\text{ В},\qquad I=0{,}100\text{ А}
    \]
    \[
        R=\frac{U}{I}=\frac{5}{0{,}100}=50\text{ Ом}
    \]
    Это оценка сопротивления в одном режиме.
    Одна точка не подтверждает линейность $I(U)$.
\end{frame}

\begin{frame}{Стандартная неопределенность}
    Для выбранной точки заданы независимые входы:
    \[
        u_U=0{,}05\text{ В},\qquad u_I=0{,}002\text{ А}.
    \]
    В первом приближении:
    \[
        u_R=\sqrt{\left(\frac{u_U}{I}\right)^2+
        \left(\frac{Uu_I}{I^2}\right)^2}\approx1{,}12\text{ Ом}.
    \]
    Это стандартная неопределенность, а не интервал с вероятностью 95\%.
\end{frame}

\begin{frame}{Вывод и ограничения}
    \begin{itemize}
        \item В выбранной точке получено $R=50$ Ом при $u_R\approx1{,}12$ Ом.
        \item Условия оценки: независимые входы и принятая модель цепи.
        \item Линейность всей серии и влияние нагрева здесь не оценивались.
    \end{itemize}
\end{frame}

\begin{frame}{Вопросы для обсуждения}
    \begin{enumerate}
        \item Как оценено влияние приборов на измеряемую цепь?
        \item Как проверить линейность по серии точек?
        \item Какие измерения помогут оценить влияние нагрева?
    \end{enumerate}
\end{frame}
\end{document}
